\documentclass[11pt]{article}

\usepackage[margin=25mm]{geometry}

\usepackage{amsmath,amssymb,fontspec,unicode-math}

\IfFontExistsTF{DejaVu Serif}{\setmainfont{DejaVu Serif}}{\setmainfont{Latin Modern Roman}}

\IfFontExistsTF{Latin Modern Math}{\setmathfont{Latin Modern Math}}{\setmathfont{DejaVu Math TeX Gyre}}

\usepackage{xurl,hyperref}

\setlength{\parindent}{0pt}\setlength{\parskip}{6pt}

\begin{document}

\title{Emergent Risk Awareness in Rational Agents under Resource Constraints}

\author{ELVIA preliminary AI analysis}\date{21 October 2026}\maketitle

Preliminary assessment: Minor.\par

Independent scientific review remains pending. Author adoption has not been requested.

The long-term-survival proof confuses a one-step Bellman expectation with realized trajectory return. An exact two-decision example disproves the required decomposition; the theorem’s complete positive-gap statement and empirical behavior are not thereby refuted.

\section*{Review scope}

Rapid, claim-focused reading of the latest applicable author-paper and resource passages. Direct mathematical or static-source reasoning is stated below. No author program, retained verifier, model inference or benchmark was executed.

\section*{Problem 1: A Bellman conditional expectation is treated as a realized return on eventual ruin}

Source: Proceedings Eq.(1), p.3; Eq.(2), p.4; Theorem 2 and preceding definitions, p.7; notation table, p.20; Appendix C, Eq.(17), p.24

Assessment: supported; impact: minor

\subsection*{Source claim (paraphrase)}

The paper defines \(G=\widetilde R+v^*_{t+1}(b+\widetilde R)\) and \(\nu=\mathbb E[G\mid S]\), with \(S\) eventual survival. Theorem 2’s proof uses \[q=-b\Pr(S^c)+\mathbb E[G\mathbf1_S].\]

\subsection*{Analysis}

Take initial budget \(b=1\), two remaining decisions, one action, and independent raw rewards \(-2,0,3\), each with probability \(1/3\). Apply the paper’s clipping \(\widetilde R=\max(-b,R)\) and absorption at zero. The one-decision continuation values are \(v(0)=0,v(1)=2/3,v(4)=1/3\). Consequently \(G\) equals \(-1,2/3,10/3\), respectively, and \(q=\mathbb EG=1\).

The nine equally likely raw pairs have realized total clipped returns \[(-1,-1,-1,-1,0,3,1,3,6)\]. Exactly the final five survive. Therefore \[\Pr(S)=5/9,\qquad\mathbb E[G\mathbf1_S]=34/27,\qquad\nu=34/15.\] The proof’s expression is \(-4/9+34/27=22/27\), differing from \(q=1\) by \(5/27\). The failure term is actually \(\mathbb E[G\mathbf1_{S^c}]=-7/27\), rather than \(-12/27\).

The path with first reward \(0\) and second reward \(-2\) explains the distinction: realized return is \(-1\), but the first-step Bellman variable remains \(2/3\). Conditioning on later ruin changes the future distribution; the unconditioned tower property does not justify the substitution. The notation table confirms the expected-continuation definition, so treating \(G\) as a realized trajectory return changes the object.

The example meets the finite, bounded reward and optimal-continuation model. With one action, however, it does not satisfy the nontrivial positive competing-action survival-gap premises of the full theorem. The theorem is not declared false by this identity counterexample. One-decision, certain-survival and immediate-ruin controls preserve the decomposition.

\subsection*{Finding and scope}

The supplied proof of the long-term-survival result needs a corrected return definition or an additional conditioning argument. The unchanged sufficient-condition theorem is not established by this proof. The other behavior results, Bellman recursion and empirical dynamic-programming plots are not refuted.

\subsection*{Proposed correction}

The supplied replacement uses realized total clipped return \(H\), which telescopes to \(-b\) on ruin. Writing \(S_a\) for survival probability and \(M_a=\mathbb E[H\mathbf1_{S_a}]\), the exact identity is \(q_a=-b(1-S_a)+M_a\). Hence \[b(S_{\hat a}-S_a)\geq M_a-M_{\hat a}\ \Longrightarrow\ q_{\hat a}\geq q_a.\] This proves a comparison with a different return-gap definition, rather than the original theorem unchanged.

The existing correction is retained as a bounded proposal; author adoption is not established.

\section*{Limits}

Preliminary AI source/argument assessment; independent human scientific adjudication remains pending.\par

Mathematical counterexamples and source discrepancies do not remeasure published experiments; historical execution results were not rerun.\par

\section*{Source identity}

NeurIPS 2025 official proceedings PDF; current endpoint freshly retrieved and relevant current author revisions compared

\url{https://papers.nips.cc/paper_files/paper/2025/file/ad52da7034c78e61a7dc1e09f1a24585-Paper-Conference.pdf}

Reviewed PDF SHA-256: \nolinkurl{0ba23df976458d8981b18951229fb924e63c8527356bffee0552af9a99357ab2}

The exact current PDF was freshly retrieved from the official endpoint. Its SHA-256 identifies the bytes inspected; the supplied historical LaTeX report remains unchanged outside this export. Paper authors were checked against the primary PDF header.

\end{document}
